\documentclass[10pt,a4paper]{amsart}
\usepackage{amsmath,amssymb,amsthm,graphicx}
% For dealing with references we use the comment environment
\usepackage{verbatim}
\newenvironment{reference}{\comment}{\endcomment}
%\newenvironment{reference}{}{}
\newenvironment{slogan}{\comment}{\endcomment}
\newenvironment{history}{\comment}{\endcomment}

% For commutative diagrams we use Xy-pic
\usepackage[all]{xy}
% Theorem environments.
%
\theoremstyle{plain}
\newtheorem{theorem}[subsection]{Theorem}
\newtheorem{proposition}[subsection]{Proposition}
\newtheorem{lemma}[subsection]{Lemma}

\theoremstyle{definition}
\newtheorem{definition}[subsection]{Definition}
\newtheorem{example}[subsection]{Example}
\newtheorem{exercise}[subsection]{Exercise}
\newtheorem{situation}[subsection]{Situation}

\theoremstyle{remark}
\newtheorem{remark}[subsection]{Remark}
\newtheorem{remarks}[subsection]{Remarks}

\numberwithin{equation}{subsection}

% Macros
%
\def\lim{\mathop{\mathrm{lim}}\nolimits}
\def\colim{\mathop{\mathrm{colim}}\nolimits}
\def\Spec{\mathop{\mathrm{Spec}}}
\def\Hom{\mathop{\mathrm{Hom}}\nolimits}
\def\Ext{\mathop{\mathrm{Ext}}\nolimits}
\def\SheafHom{\mathop{\mathcal{H}\!\mathit{om}}\nolimits}
\def\SheafExt{\mathop{\mathcal{E}\!\mathit{xt}}\nolimits}
\def\Sch{\mathit{Sch}}
\def\Mor{\mathop{\mathrm{Mor}}\nolimits}
\def\Ob{\mathop{\mathrm{Ob}}\nolimits}
\def\Sh{\mathop{\mathit{Sh}}\nolimits}
\def\NL{\mathop{N\!L}\nolimits}
\def\CH{\mathop{\mathrm{CH}}\nolimits}
\def\proetale{{pro\text{-}\acute{e}tale}}
\def\etale{{\acute{e}tale}}
\def\QCoh{\mathit{QCoh}}
\def\Ker{\mathop{\mathrm{Ker}}}
\def\Im{\mathop{\mathrm{Im}}}
\def\Coker{\mathop{\mathrm{Coker}}}
\def\Coim{\mathop{\mathrm{Coim}}}

% Boxtimes
%
\DeclareMathSymbol{\boxtimes}{\mathbin}{AMSa}{"02}

%
% Macros for moduli stacks/spaces
%
\def\QCohstack{\mathcal{QC}\!\mathit{oh}}
\def\Cohstack{\mathcal{C}\!\mathit{oh}}
\usepackage{cleveref}
\crefname{originalsection}{original section}{original sections}
\Crefname{originalsection}{Original section}{Original sections}
\makeatletter
\newlabel{derived-remark-map-into-derived-limit-truncations}{{0H72}{source}{}{sourcepointer-1}{}}
\newlabel{derived-remark-map-into-derived-limit-truncations@cref}{{[originalsection][0][]0H72}{[1][0][]source}{}{}{}}
\makeatother
\crefname{originalsection}{original section}{original sections}
\Crefname{originalsection}{Original section}{Original sections}
\makeatletter
\newlabel{derived-remark-truncation-distinguished-triangle}{{08J5}{source}{}{sourcepointer-2}{}}
\newlabel{derived-remark-truncation-distinguished-triangle@cref}{{[originalsection][0][]08J5}{[1][0][]source}{}{}{}}
\makeatother
\crefname{originallemma}{original lemma}{original lemmas}
\Crefname{originallemma}{Original lemma}{Original lemmas}
\makeatletter
\newlabel{lemma-RGamma-commutes-with-Rlim}{{0D60}{source}{}{sourcepointer-3}{}}
\newlabel{lemma-RGamma-commutes-with-Rlim@cref}{{[originallemma][0][]0D60}{[1][0][]source}{}{}{}}
\makeatother
\crefname{originallemma}{original lemma}{original lemmas}
\Crefname{originallemma}{Original lemma}{Original lemmas}
\makeatletter
\newlabel{lemma-sheafification-cohomology}{{0BKJ}{source}{}{sourcepointer-4}{}}
\newlabel{lemma-sheafification-cohomology@cref}{{[originallemma][0][]0BKJ}{[1][0][]source}{}{}{}}
\makeatother
\title{Pointwise diagonal cohomological vanishing implies completeness of the Postnikov tower}
\author{The Stacks Project Authors}
\date{stacks-project:cohomology.tex:local-snapshot-20261011; GNU FDL 1.2 or later}
\begin{document}
\maketitle
\noindent\textit{Editorial note.} The bound n(x) is chosen separately for each fixed cohomological degree m; no global bound on cohomological dimension is assumed.\par\medskip
\noindent\textit{Editorial note.} The stated canonical comparison is understood in the standard compatible-map sense of the cited derived-limit remark; the proof shows every compatible choice is an isomorphism.\par\medskip
\noindent\textit{Editorial scope.} One original target; complete source-local proof chain included, standard derived-sheaf primitives explicitly delimited.\par\medskip
\noindent\textit{Allowed standard prerequisites (editorial).}\begin{itemize}
\item Derived inverse limits and Milnor sequence. RΓ(U,-) commutes with derived countable inverse limits; its cohomology lies in the short exact sequence with lim and R\textasciicircum{}1lim. Eventually constant inverse systems have R\textasciicircum{}1lim=0.
\item Standard truncation triangles and sheafification. Canonical truncations fit in the cohomology-object triangle, and H\textasciicircum{}m(K) is the sheaf associated to U\ensuremath{\mapsto}H\textasciicircum{}m(U,K).
\end{itemize}\medskip
\ifdefined\hypertarget\hypertarget{sourcepointer-1}{}\fi
\noindent\textit{Source reference.} Original section 0H72: focused-seeds/ACL/sources/stacks-snapshot-20261011/derived.tex, source lines 10557--10586. Standard definition of the compatible map to the truncation inverse limit. Displayed pointer is the original stable Stacks tag.\par\medskip
\ifdefined\hypertarget\hypertarget{sourcepointer-2}{}\fi
\noindent\textit{Source reference.} Original section 08J5: focused-seeds/ACL/sources/stacks-snapshot-20261011/derived.tex, source lines 4075--4109. Standard truncation triangle; it does not establish the pointwise vanishing criterion. Displayed pointer is the original stable Stacks tag.\par\medskip
\ifdefined\hypertarget\hypertarget{sourcepointer-3}{}\fi
\noindent\textit{Source reference.} Original lemma 0D60: focused-seeds/ACL/sources/stacks-snapshot-20261011/cohomology.tex, source lines 9456--9467. Explicit permitted Milnor exact sequence. Displayed pointer is the original stable Stacks tag.\par\medskip
\ifdefined\hypertarget\hypertarget{sourcepointer-4}{}\fi
\noindent\textit{Source reference.} Original lemma 0BKJ: focused-seeds/ACL/sources/stacks-snapshot-20261011/cohomology.tex, source lines 7875--7883. Standard description of cohomology sheaves, used only to choose and shrink representatives. Displayed pointer is the original stable Stacks tag.\par\medskip
\begin{lemma}
\label{lemma-is-limit-per-point}
Let $(X, \mathcal{O}_X)$ be a ringed space. Let $E \in D(\mathcal{O}_X)$.
Assume that for every $x \in X$ there exist
a function $p(x, -) : \mathbf{Z} \to \mathbf{Z}$ and
a fundamental system $\mathfrak{U}_x$ of open neighbourhoods of $x$
such that
$$
H^p(U, H^{m - p}(E)) = 0 \text{ for }
U \in \mathfrak{U}_x \text{ and } p > p(x, m)
$$
Then the map $E \to R\lim \tau_{\geq -n} E$ of
Derived Categories, Remark
\ref{derived-remark-map-into-derived-limit-truncations}
is an isomorphism in $D(\mathcal{O}_X)$.
\end{lemma}

\begin{lemma}
\label{lemma-cohomology-derived-limit-injective}
Let $(X, \mathcal{O}_X)$ be a ringed space. Let $(K_n)$ be an
inverse system in $D(\mathcal{O}_X)$. Let $x \in X$ and $m \in \mathbf{Z}$.
Assume there exist an integer $n(x)$ and a fundamental system $\mathfrak{U}_x$
of open neighbourhoods of $x$ such that for $U \in \mathfrak{U}_x$
\begin{enumerate}
\item $R^1\lim H^{m - 1}(U, K_n) = 0$, and
\item $H^m(U, K_n) \to H^m(U, K_{n(x)})$ is injective
for $n \geq n(x)$.
\end{enumerate}
Then the map on stalks $H^m(R\lim K_n)_x \to H^m(K_{n(x)})_x$ is injective.
\end{lemma}

\begin{proof}
Let $\gamma$ be an element of $H^m(R\lim K_n)_x$ which maps to zero
in $H^m(K_{n(x)})_x$. Since $H^m(R\lim K_n)$ is the sheafification
of $U \mapsto H^m(U, R\lim K_n)$
(by Lemma \ref{lemma-sheafification-cohomology})
we can choose $U \in \mathfrak{U}_x$
and an element $\tilde \gamma \in H^m(U, R\lim K_n)$ mapping to $\gamma$.
Then $\tilde\gamma$ maps to $\tilde\gamma_{n(x)} \in H^m(U, K_{n(x)})$.
Using that $H^m(K_{n(x)})$ is the sheafification of
$U \mapsto H^m(U, K_{n(x)})$
(by Lemma \ref{lemma-sheafification-cohomology} again)
we see that after shrinking $U$ we may assume that $\tilde\gamma_{n(x)} = 0$.
For this $U$ we consider the short exact sequence
$$
0 \to
R^1\lim H^{m - 1}(U, K_n) \to H^m(U, R\lim K_n) \to
\lim H^m(U, K_n) \to 0
$$
of Lemma \ref{lemma-RGamma-commutes-with-Rlim}.
By assumption (1) the group on the left is zero and by
assumption (2) the group on the right maps injectively
into $H^m(U, K_{n(x)})$. We conclude $\tilde\gamma = 0$
and hence $\gamma = 0$ as desired.
\end{proof}

\begin{proof}
Set $K_n = \tau_{\geq -n}E$ and $K = R\lim K_n$.
The canonical map $E \to K$
comes from the canonical maps $E \to K_n = \tau_{\geq -n}E$.
We have to show that $E \to K$ induces an isomorphism
$H^m(E) \to H^m(K)$ of cohomology sheaves. In the rest of the
proof we fix $m$. If $n \geq -m$, then
the map $E \to \tau_{\geq -n}E = K_n$ induces an isomorphism
$H^m(E) \to H^m(K_n)$.
To finish the proof it suffices to show that for every $x \in X$
there exists an integer $n(x) \geq -m$ such that the map
$H^m(K)_x \to H^m(K_{n(x)})_x$ is injective. Namely, then
the composition
$$
H^m(E)_x \to H^m(K)_x \to H^m(K_{n(x)})_x
$$
is a bijection and the second arrow is injective, hence the
first arrow is bijective. Set
$$
n(x) = 1 + \max\{-m, p(x, m - 1) - m, -1 + p(x, m) - m, -2 + p(x, m + 1) - m\}.
$$
so that in any case $n(x) \geq -m$. Claim: the maps
$$
H^{m - 1}(U, K_{n + 1}) \to H^{m - 1}(U, K_n)
\quad\text{and}\quad
H^m(U, K_{n + 1}) \to H^m(U, K_n)
$$
are isomorphisms for $n \geq n(x)$ and $U \in \mathfrak{U}_x$.
The claim implies conditions
(1) and (2) of Lemma \ref{lemma-cohomology-derived-limit-injective}
are satisfied and hence implies the desired injectivity.
Recall (Derived Categories, Remark
\ref{derived-remark-truncation-distinguished-triangle})
that we have distinguished triangles
$$
H^{-n - 1}(E)[n + 1] \to
K_{n + 1} \to K_n \to H^{-n - 1}(E)[n + 2]
$$
Looking at the associated long exact cohomology sequence the claim follows if
$$
H^{m + n}(U, H^{-n - 1}(E)),\quad
H^{m + n + 1}(U, H^{-n - 1}(E)),\quad
H^{m + n + 2}(U, H^{-n - 1}(E))
$$
are zero for $n \geq n(x)$ and $U \in \mathfrak{U}_x$.
This follows from our choice of $n(x)$
and the assumption in the lemma.
\end{proof}

\end{document}
